% !TEX root = diss.tex

\chapter{Discussion and Future Directions}

In this work, we extended optimal transport trajectory inference for single-cell data in two key facets. First, we proposed a methodology for constructing trajectory-informed features from collections of single-timepoint snapshots, using OT as the underlying temporal kernel (Chapter 2). Second, we established that OT is capable of modeling population dynamics of CD8 T cells across tissues during acute viral infection and predicted distinct fates for tissue-resident CD8 T cells based on time of arrival in peripheral tissues (Chapter 3).

\section{Featurization of Single Cell Trajectories Through Kernel Mean Embedding of Optimal Transport Maps}

In chapter 2, we introduced the problem of "featurizing"---or creating computational summaries of---single cell trajectories given snapshot data. OT and related methods are capable of efficiently and accurately fitting trajectories to time-resolved data \citep{schiebingerOptimalTransportAnalysisSingleCell2019, kleinMappingCellsTime2025, kleinCellFlowEnablesGenerative2025}. However the transport plan, the kernel of OT, is inherently high-dimensional, being composed of a dense matrix for just a single pair of timepoints. In analyses involving more than a handful of timepoints, this dimensionality problem hinders our ability to perform cell-based trajectory analysis tasks such as clustering or differential abundance testing over the entire experimental time span. The featurization approach introduced in this chapter addresses this shortcoming by projecting cells across the entirety of a trajectory into a shared feature space.

We derived that transforming features using a reproducing kernel function, followed by OT projection through time, produces a trajectory kernel mean embedding (TKME) of cell trajectories. In practice, this means that it is possible to encode not only the mean state of a cell over time, but also higher order moments, capturing uncertainty arising from trajectory branching or cell death. We benchmarked this approach for downstream clustering of HSC differentiation, finding that clusters fit on TKME produced more consistent partitions of the OT trajectories than those fit on OT projections into PCA space or those fit without using any trajectory information.

Using TKME embeddings allowed us to leverage kernel-based two sample tests to quantify the effects of experimental perturbations on cell differentiation. This approach measured not just static distributional shifts in cell populations due to perturbations, but also the temporal dynamics of these shifts, providing a robust test for experimentally impactful perturbations. We demonstrated that this approach strengthened detection of transcription factors with functional roles in iPSC differentiation, with our approach elevating two additional transcription factors (\textit{MYC}, \textit{NANOG}) having confirmed biological activity \citep{ishikawaRENGEInfersGene2023} above our detection threshold. 

Future work on this area should focus on developing more efficient compressed trajectory representations. Currently, TKME scales with the number of sampled timepoints, which is inefficient for finely-sampled trajectories. One potential solution is to learn a trajectory representation parameterized by a deep neural network. Similarly, another extension is to generate new trajectories by sampling a latent space. Recent works have begun exploring this area \citep{xuVISTVariationalInference2026}, however there is still a need for methods that explicitly incorporate population dynamics and branching trajectories.

\section{Optimal Transport Fate Mapping Resolves T Cell Differentiation Dynamics Across Tissues}

In chapter 3, we applied OT trajectory inference to the CD8 T cell response to acute infection, integrating TKME with single cell clustering to develop a "fate flow" framework for trajectory analysis.

There are several aspects of the T cell infection response that make OT a natural choice for trajectory modeling. For example, the timescale of the infection response is on the order of days to months, so multiple timepoints are necessary to profile the entire trajectory. Furthermore, CD8 T cell states during the infection response are cyclic, starting with all cells in a stem-like naive state, transiting through terminally-differentiated populations at the peak of infection, before returning to a stem-like memory population after antigen clearance. This cyclicity has been observed to misdirect other trajectory inference approaches \citep{hudsonTechnologyMeetsTILs2022}, yet the temporal ordering imposed by OT inherently accounts for it. Finally, the outsize role of proliferation during clonal expansion and apoptosis during contraction can confound trajectories. Therefore, we adopted an OT trajectory framework that enabled us to proactively estimate population growth rates from proliferation and apoptosis gene signatures, enabling us to explicitly control for the confounding effects of population dynamics.

We first tested the ability of OT to model CD8 T cell dynamics in circulation. We found that estimated growth rates recapitulated population dynamics of clonal expansion and contraction, and that trajectories partitioned into flows based on memory and effector gene signatures. Subsetting the CD8 T cells also produced coherent flows, suggesting that OT trajectories agreed with prior knowledge about CD8 T cell differentiation. Clustering based on TKME simplified the major trajectories, revealing stable lineages corresponding to short-lived early effector, terminal effector, and long-lived central memory subsets. Interestingly, this analysis suggested that there is a larger basin for memory-fated T cells than expected, which implies that cells with effector characteristics may be able to de-differentiate into central memory cells. This fits with some emerging theories about flexible fate decisions in CD8 T cell memory formation \citep{abadieReversibleTunableEpigenetic2024}, however further validation is needed to assess whether these effector-like memory precursors can be detected in other datasets.

Next, we adapted OT to a novel setting, modeling the migration of CD8 T cells between biological compartments. We found that this model recapitulated the experimentally determined timeline of migration from circulation into the siIEL compartment \citep{masopustDynamicCellMigration2010}, and was able to correctly identify early and late arrivals to the siIEL, which we tested using an intravenous antibody labeling platform. A more challenging problem is identifying whether there are some circulating T cell populations with inherent propensities for migration into peripheral tissues. Our approach predicted that while T\textsc{RM} may originate from diverse lineages, memory-precursor lineages had the greatest propensity for tissue migration. This fits with reported phenotypes of T\textsc{RM}-biased clones in previous studies \citep{kokPrecursorsCD8Tissue2022,kokCommittedTissueresidentMemory2020}. Further, we predicted that the migration propensities of different lineages change over time, with true T\textsc{RM} precursors being more likely to migrate into tissues early, while more transient, short-lived effectors make up the majority of late arrivals.

Finally, we compared differentiation trajectories of circulating and resident memory CD8 T cells once they are established within each compartment, using TKME clusters to facilitate balanced comparisons between trajectories over time. This revealed shared and distinct regulatory patterns within both circulating and tissue resident memory. Short-lived effector signatures were highly conserved between both compartments, as were effector memory signatures. However, the most stem-like T\textsc{CM} lineages and progenitor-like T\textsc{RM} lineages showed substantial divergence in their gene signatures. This indicated context-dependent roles for transcription factors during CD8 T cell differentiation. For example, \textit{Tbx21} was enriched in effector trajectories in both tissues, while \textit{Tcf7} showed a much higher degree of enrichment in circulating memory than resident memory trajectories. We also nominated \textit{Tfap4} as an early regulator of circulating effector memory. Together, these results reveal important distinctions between compartment-specific CD8 T cell differentiation.

Future work should seek to extend CD8 T cell trajectory reconstruction to new systems, for example other infections or cancer, while further strengthening the modeling approach. For example, OT trajectories can be regularized by incorporating lineage information such as clonal identity inferred from TCR sequencing data of endogenous infection models. In addition, incorporating intermediate compartments (such as draining lymph nodes) and higher resolution timepoints would improve the fidelity of the model system. Computationally, using sequential OT is prone to create blurry mappings at long ranges. This can be controlled via multi-marginal OT, which solves OT mappings for the entire trajectory simultaneously. Using OT to train generative trajectory models, e.g. neural ODEs using OT-conditional flow matching \citep{tongTrajectoryNetDynamicOptimal2020, tongImprovingGeneralizingFlowbased2024, tongSimulationfreeSchrodingerBridges2024}, would enable simulating experimental perturbations such as transcription factor knock out or over-expression.
